\cleardoublepage
\phantomsection
\chapter*{校译札记}
\addcontentsline{toc}{chapter}{校译札记}
\markboth{校译札记}{校译札记}

% 请在此处填写“校译札记”正文。

Morse理论在流形拓扑的研究中占有核心地位，而$h$配边定理以极其生动的方式向我们展示了Morse理论是如何实现这一点的。拓扑学大师John Milnor的三部曲《从微分观点看拓扑》，《Morse理论》和《$h$配边定理讲义》则用精湛而优雅的语言向我们展示了其中的精要之处。

时至今日，《从微分观点看拓扑》和《Morse理论》已有翻译版问世，前者由熊金诚教授翻译，而后者由江嘉禾教授翻译。而《$h$配边定理讲义》，自初识此书，校者便热切期盼它能以中文的形式传播。尽管对于众多数学学习者，阅读英文已不是难事。然而对于相当一部分读者而言，阅读中文仍更为方便。而且，从中文文献的完整性来说，如此阙如令校者颇为遗憾。为了弥补这个缺憾，校者从2020年便启动了这个翻译计划。然而，学业和工作任务繁重，加上校者的懒惰，六年过去，手工译稿仅有3页完成。

令人颇为感慨的是，六年后的今天大语言模型的发展日新月异。时而有如核爆瘫坐。不免让校者日夜思索：人类社群应该如何面对知识的诞生和存留。诚然，校者尚未有明确答案。只是扪心自问，在时代的洪流中，我要如何将我作为人类知识的参与者的痕迹保留下一些？左思右想之后，我想在人人皆抽卡的时代，我还是希望以一种冷静的态度面对世界的变化，为此一个简单的回应便是当前的译稿，使之作为我合理使用了一些词元的证明。

尽管译稿仍由大语言模型自主翻译完成，校者仍然费了些许心力在校对和制图上。尤其是囿于大语言模型识图制图能力上的欠缺，本译本中的大多数图像仍由校者经由\hyperlink{https://www.mathcha.io/}{Mathcha}绘制。在内容上，校者选择尽可能在文字上还原原著的风貌，乃至保留了若干原著中无关紧要的排版错误；而若干中文翻译尚未得到共识，作者的私心是传播一些自己喜欢的翻译。此外，校者也参考了\hyperlink{https://oldbookstonew.blogspot.com/}{Old books with high-quality typesetting}项目
的高质量 \LaTeX 重印本。

当然，所有尚未被校出的翻译错误，本人均为此负全责。也欢迎各位路人、路机指出问题。


\vfill
\begin{flushright}
校者：张秉宇\par
\vspace{1em}
日期：2026年8月26日\par
\end{flushright}
